\documentclass[UTF8]{ctexart}
\usepackage{amsmath, amssymb, geometry}
\usepackage{xcolor}
\usepackage{enumitem}
\usepackage{tcolorbox}
\usepackage{cases}
\usepackage{fancyhdr}
\geometry{a4paper, margin=1in}

% 定义红色环境
\newenvironment{redenv}{\color{red}}{}

% 定义详细解析环境
\newenvironment{detailedanalysis}{%
    \color{red}
    \begin{enumerate}[label=\textbf{第\arabic*步：}, leftmargin=*, align=left, itemsep=1.5ex]
}{%
    \end{enumerate}
}

% 定义概念框
\newenvironment{conceptbox}{%
    \begin{tcolorbox}[colback=red!5, colframe=red!50!black, title=概念解析, fonttitle=\bfseries]
}{%
    \end{tcolorbox}
}

% 定义公式推导环境
\newenvironment{formuladerivation}{%
    \begin{enumerate}[label={\textbf{公式\arabic*：}}, leftmargin=*, align=left]
}{%
    \end{enumerate}
}

% 定义题目和解析的分隔线
\newcommand{\problemrule}{\noindent\rule{\textwidth}{1.5pt}\vspace{-0.8em}}
\newcommand{\analysisrule}{\noindent\rule{\textwidth}{0.8pt}\vspace{-0.8em}}

% 宏定义：方便排版选择题选项
\newcommand{\choicesFour}[4]{
    \begin{enumerate}[label=\Alph*., itemsep=0pt, parsep=0pt, topsep=5pt, labelsep=0.5em]
        \begin{minipage}{0.22\linewidth} \item #1 \end{minipage}
        \begin{minipage}{0.22\linewidth} \item #2 \end{minipage}
        \begin{minipage}{0.22\linewidth} \item #3 \end{minipage}
        \begin{minipage}{0.22\linewidth} \item #4 \end{minipage}
    \end{enumerate}
}
\newcommand{\choicesTwo}[4]{
    \begin{enumerate}[label=\Alph*., itemsep=0pt, parsep=0pt, topsep=5pt, labelsep=0.5em]
        \begin{minipage}{0.45\linewidth} \item #1 \end{minipage}
        \begin{minipage}{0.45\linewidth} \item #2 \end{minipage}
        \begin{minipage}{0.45\linewidth} \item #3 \end{minipage}
        \begin{minipage}{0.45\linewidth} \item #4 \end{minipage}
    \end{enumerate}
}
\newcommand{\choices}[4]{
    \begin{enumerate}[label=\Alph*., itemsep=0pt, parsep=0pt, topsep=5pt, labelsep=0.5em]
        \item #1
        \item #2
        \item #3
        \item #4
    \end{enumerate}
}

\newcommand{\blank}[1]{\underline{\hspace{#1}}}

\begin{document}

\section*{第十一章 二重积分（提高）}

% ========== 第1题 ==========
\problemrule
\section*{【题目1】交换积分次序}
设函数 $ f(x,y) $ 连续，则 $ \int_{1}^{2}dy\int_{y}^{2}f(x,y)dx $ 交换积分次序为 \blank{1cm}。
\choicesTwo
{$ \int_{1}^{2}dx\int_{1}^{x}f(x,y)dy $}
{$ \int_{1}^{2}dx\int_{x}^{1}f(x,y)dy $}
{$ \int_{1}^{2}dx\int_{2}^{x}f(x,y)dy $}
{$ \int_{1}^{2}dx\int_{x}^{2}f(x,y)dy $}

\begin{redenv}
\analysisrule
\subsection*{逐步解析}

\begin{detailedanalysis}
\item \textbf{详细求解过程}
\begin{enumerate}[label=(\arabic*)]
\item \textbf{确定积分区域 D}：
原序：$1 \le y \le 2$，$y \le x \le 2$。
画图可知：区域由 $y=1, x=2, y=x$ 围成。
\item \textbf{交换次序}：
先积 $y$，后积 $x$。
$x$ 的范围是从 $1$ 到 $2$。
对固定的 $x$， $y$ 从 $1$ 到 $x$。
$1 \le x \le 2$，$1 \le y \le x$。
\item \textbf{结果}：
$\int_1^2 dx \int_1^x f(x,y) dy$。
\end{enumerate}

\item \textbf{最终答案}
\textbf{答案}：A
\end{detailedanalysis}
\end{redenv}

% ========== 第2题 ==========
\problemrule
\section*{【题目2】方形域二重积分}
设 $D$ 是方形域： $ 0 \leq x \leq 1, 0 \leq y \leq 1, \iint_{D} xyd\sigma = $ \blank{1cm}。
\choicesFour{1}{$ \frac{1}{2} $}{$ \frac{1}{3} $}{$ \frac{1}{4} $}

\begin{redenv}
\analysisrule
\subsection*{逐步解析}

\begin{detailedanalysis}
\item \textbf{详细求解过程}
可分离变量：
$$ \iint_D xy dxdy = (\int_0^1 x dx) (\int_0^1 y dy) = [\frac{1}{2}x^2]_0^1 \cdot [\frac{1}{2}y^2]_0^1 = \frac{1}{2} \cdot \frac{1}{2} = \frac{1}{4} $$

\item \textbf{最终答案}
\textbf{答案}：D
\end{detailedanalysis}
\end{redenv}

% ========== 第3题 ==========
\problemrule
\section*{【题目3】对称性与交换次序}
$ I = \int_{0}^{1} dy \int_{-\sqrt{1-y}}^{\sqrt{1-y}} 3x^{2}y^{2} dx $ 交换积分次序为 $I =$ \blank{1cm}。
\choicesTwo
{$ \int_{0}^{1}dy\int_{0}^{\sqrt{1-y}}3x^{2}y^{2}dx $}
{$ \int_{0}^{\sqrt{1-y}}dy\int_{0}^{1}3x^{2}y^{2}dx $}
{$ \int_{-1}^{1}dx\int_{0}^{1-x^{2}}3x^{2}y^{2}dy $}
{$ \int_{0}^{1}dx\int_{0}^{1+x^{2}}3x^{2}y^{2}dy $}

\begin{redenv}
\analysisrule
\subsection*{逐步解析}

\begin{detailedanalysis}
\item \textbf{详细求解过程}
\begin{enumerate}[label=(\arabic*)]
\item \textbf{确定区域}：
$0 \le y \le 1$，$ -\sqrt{1-y} \le x \le \sqrt{1-y} $。
即 $x^2 \le 1-y \implies y \le 1-x^2$。且 $y \ge 0$。
$x$范围：$-\sqrt{1} \le x \le \sqrt{1} \implies -1 \le x \le 1$。
\item \textbf{交换次序}：
先积 $y$，后积 $x$。
$-1 \le x \le 1$，$0 \le y \le 1-x^2$。
\item \textbf{对应选项}：
选项 C 符合形式。
\end{enumerate}

\item \textbf{最终答案}
\textbf{答案}：C
\end{detailedanalysis}
\end{redenv}

% ========== 第4题 ==========
\problemrule
\section*{【题目4】极坐标面积元素}
二重积分 $ \iint_{D} f(x,y)d\sigma $ 在极坐标下的面积元素为 \blank{1cm}。
\choicesFour{$ d\sigma = dx dy $}{$ d\sigma = r dr d\theta $}{$ d\sigma = dr d\theta $}{$ d\sigma = r^{2}\sin\theta dr d\theta $}

\begin{redenv}
\analysisrule
\subsection*{逐步解析}

\begin{detailedanalysis}
\item \textbf{详细求解过程}
极坐标变换 $x = r\cos\theta, y = r\sin\theta$。
雅可比行列式 $J = r$。
面积元素 $d\sigma = |J| dr d\theta = r dr d\theta$。

\item \textbf{最终答案}
\textbf{答案}：B
\end{detailedanalysis}
\end{redenv}

% ========== 第5题 ==========
\problemrule
\section*{【题目5】交换积分顺序}
设 $ I = \int_{0}^{1} dx \int_{x^{2}}^{x} f(x, y) dy $ 交换积分顺序后，则 $ I = $ \blank{2cm}。

\begin{redenv}
\analysisrule
\subsection*{逐步解析}

\begin{detailedanalysis}
\item \textbf{详细求解过程}
\begin{enumerate}[label=(\arabic*)]
\item \textbf{确定区域}：
$0 \le x \le 1$，$x^2 \le y \le x$。
由 $y=x^2$ 和 $y=x$ 围成。
\item \textbf{交换次序}：
$y$ 的范围也是 $0 \le y \le 1$。
对固定 $y$，左边界 $x=y$，右边界 $x=\sqrt{y}$。
所以 $y \le x \le \sqrt{y}$。
\item \textbf{结果}：
$\int_0^1 dy \int_y^{\sqrt{y}} f(x,y) dx$。
\end{enumerate}

\item \textbf{最终答案}
\textbf{答案}：$ \int_0^1 dy \int_y^{\sqrt{y}} f(x,y) dx $
\end{detailedanalysis}
\end{redenv}

% ========== 第6题 ==========
\problemrule
\section*{【题目6】极坐标积分}
$D$ 是圆环 $ a^{2}\leq x^{2}+y^{2}\leq b^{2} $ ，计算 $ I=\iint_{D}(x^{2}+y^{2})d\sigma $。

\begin{redenv}
\analysisrule
\subsection*{逐步解析}

\begin{detailedanalysis}
\item \textbf{详细求解过程}
使用极坐标：
$x^2+y^2 = r^2$。
区域 $a \le r \le b, 0 \le \theta \le 2\pi$。
$$ I = \int_0^{2\pi} d\theta \int_a^b r^2 \cdot r dr = 2\pi \int_a^b r^3 dr = 2\pi [\frac{1}{4}r^4]_a^b = \frac{\pi}{2}(b^4 - a^4) $$

\item \textbf{最终答案}
\textbf{答案}：$ \frac{\pi}{2}(b^4 - a^4) $
\end{detailedanalysis}
\end{redenv}

% ========== 第7题 ==========
\problemrule
\section*{【题目7】分段积分交换}
交换积分顺序： $ \int_{0}^{\frac{1}{4}}dy\int_{y}^{\sqrt{y}}f(x,y)dx+\int_{\frac{1}{4}}^{\frac{1}{2}}dy\int_{y}^{\frac{1}{2}}f(x,y)dx $ 为 \blank{2cm}。

\begin{redenv}
\analysisrule
\subsection*{逐步解析}

\begin{detailedanalysis}
\item \textbf{详细求解过程}
\begin{enumerate}[label=(\arabic*)]
\item \textbf{画图分析}：
Region 1: $0 \le y \le 1/4$, $y \le x \le \sqrt{y}$。
Region 2: $1/4 \le y \le 1/2$, $y \le x \le 1/2$。
\item \textbf{边界曲线}：
左边界 $x=y$。
右边界 $x=\sqrt{y} \implies y=x^2$ 和 $x=1/2$。
下边界 $y=0$。
实际上，这两个区域拼成了一个更大的区域。
观察 $x$ 的范围：
从 $R_1$ 看，当 $y=1/4$ 时，$x_{max} = \sqrt{1/4} = 1/2$。
总体 $x$ 从 $0$ 到 $1/2$。
下界 $y$：从 $x$ 轴 $y=0$ 开始？不是，$R_1$ 下界 $y=0$。
上界 $y$：
$x=y \implies y=x$。
$x=\sqrt{y} \implies y=x^2$。
整个区域是由 $y=x^2$ (下) 和 $y=x$ (上) 围成的吗？
检查 $R_1$：$y \le x \le \sqrt{y}$。即 $x$ 在直线右，抛物线左。
检查 $R_2$：$y \le x \le 1/2$。
$R_1$ 的右边界是抛物线 $x=\sqrt{y}$，即 $y=x^2$。
$R_2$ 的右边界是直线 $x=1/2$。
其实 $R_1$ 和 $R_2$ 是关于 $y=1/4$ 分割的。
在 $x \in [0, 1/2]$ 上：
下曲线是 $y=x^2$（来自 $x=\sqrt{y}$ 的反函数，注意 $y \le 1/4$ 对应 $x \le 1/2$）。
上曲线是 $y=x$。
$x$ 的最大范围是 $1/2$。
所以合并后区域为：$0 \le x \le 1/2$， $x^2 \le y \le x$。
\item \textbf{结果}：
$\int_0^{1/2} dx \int_{x^2}^x f(x,y) dy$。
\end{enumerate}

\item \textbf{最终答案}
\textbf{答案}：$ \int_0^{1/2} dx \int_{x^2}^x f(x,y) dy $
\end{detailedanalysis}
\end{redenv}

% ========== 第8题 ==========
\problemrule
\section*{【题目8】极坐标计算}
计算二重积分 $ \iint_{D} \frac{y}{x} d\sigma $ ，其中 $D$ 由 $ x^{2} + y^{2} \leq a^{2} (a > 0) $ ，$y = x$ 及 $x$ 轴在第一象限所围成的区域。

\begin{redenv}
\analysisrule
\subsection*{逐步解析}

\begin{detailedanalysis}
\item \textbf{详细求解过程}
\begin{enumerate}[label=(\arabic*)]
\item \textbf{确定区域}：
极坐标下：$y=0 \implies \theta=0$；$y=x \implies \theta=\pi/4$。
半径 $0 \le r \le a$。
\item \textbf{变换积分}：
$\frac{y}{x} = \tan\theta$。
$d\sigma = r dr d\theta$。
$$ I = \int_0^{\pi/4} d\theta \int_0^a \tan\theta \cdot r dr = \int_0^{\pi/4} \tan\theta d\theta \cdot [\frac{1}{2}r^2]_0^a $$
$$ = \frac{a^2}{2} [-\ln|\cos\theta|]_0^{\pi/4} = \frac{a^2}{2} (-\ln\frac{\sqrt{2}}{2} + \ln 1) = \frac{a^2}{2} \ln\sqrt{2} = \frac{a^2}{4} \ln 2 $$
\end{enumerate}

\item \textbf{最终答案}
\textbf{答案}：$ \frac{a^2}{4} \ln 2 $
\end{detailedanalysis}
\end{redenv}

% ========== 第9题 ==========
\problemrule
\section*{【题目9】极坐标计算}
二重积分 $ I=\iint_{D}\ln(1+x^{2}+y^{2})dxdy $ ,其中 $D$ 是由圆周 $ x^{2}+y^{2}=1 $ 及坐标轴所围成的第一象限内的区域。

\begin{redenv}
\analysisrule
\subsection*{逐步解析}

\begin{detailedanalysis}
\item \textbf{详细求解过程}
极坐标：$0 \le \theta \le \pi/2$，$0 \le r \le 1$。
$$ I = \int_0^{\pi/2} d\theta \int_0^1 \ln(1+r^2) r dr = \frac{\pi}{2} \cdot \frac{1}{2} \int_0^1 \ln(1+r^2) d(1+r^2) $$
令 $u = 1+r^2$，$u \in [1, 2]$。
$$ = \frac{\pi}{4} \int_1^2 \ln u du = \frac{\pi}{4} [u\ln u - u]_1^2 = \frac{\pi}{4} (2\ln 2 - 2 - (0 - 1)) = \frac{\pi}{4} (2\ln 2 - 1) $$

\item \textbf{最终答案}
\textbf{答案}：$ \frac{\pi}{4} (2\ln 2 - 1) $
\end{detailedanalysis}
\end{redenv}

% ========== 第10题 ==========
\problemrule
\section*{【题目10】广义变换（或直角坐标）}
求二重积分 $ \iint_{D} (x + 2y) \, dx \, dy $ ，其中区域 $D$ 是由 $ y = x, y = 4x, xy = 1 $ 围成的第一象限部分的闭区域。

\begin{redenv}
\analysisrule
\subsection*{逐步解析}

\begin{detailedanalysis}
\item \textbf{详细求解过程}
\begin{enumerate}[label=(\arabic*)]
\item \textbf{确定交点}：
$y=x, xy=1 \implies x^2=1 \implies (1,1)$。
$y=4x, xy=1 \implies 4x^2=1 \implies (1/2, 2)$。
\item \textbf{积分计算}：
分为左右两块，或先积 $x$。
固定 $y$:
从 $y=1$ 到 $y=2$：左界 $x=y/4$，右界 $x=1/y$？ No.
看图：$y=x$ 下界，$y=4x$ 上界，$xy=1$ 右界。
交点 $(1,1)$ 和 $(1/2, 2)$。
显然区域在 $x \in [1/2, 1]$ 之间。
分块：
$D_1$: $x \in [1/2, \sqrt{\dots}]$? 比较复杂。
使用换元法：令 $u=y/x, v=xy$。
$y=x \implies u=1$。$y=4x \implies u=4$。
$xy=1 \implies v=1$。
第四个边界？原题似乎隐含了 $xy=?$ 或者就是三角形区域。
题目只给了 $y=x, y=4x, xy=1$ 三条线。这是一个曲边三角形。
范围：$1 \le u \le 4$。
$v$ 的范围？
$xy=1$ 是上界。下界是原点？不，三条线围成。
交点对应 $v=1$。
原点 $0$ 显然不在。
Wait, $y=x$ 与 $y=4x$ 交于原点。
积分区域是由 $(0,0), (1,1), (1/2, 2)$ 围成？
是的。
在 $uv$ 平面上：
$u=1, u=4$。
$v$ 从 $0$ 到 $1$。
$x = \sqrt{v/u}, y = \sqrt{uv}$。
$J = \frac{\partial(x,y)}{\partial(u,v)} = \dots$
也可直接直角坐标分块：
$x$ 从 $0$ 到 $1/2$：$y$ 从 $x$ 到 $4x$。（$0\to 1/2$ 这一段被 $xy=1$ 切断吗？ $x=1/2$ 时 $y=2$，正好是 $y=4x$ 与 $xy=1$ 交点）。
$x$ 从 $1/2$ 到 $1$：$y$ 从 $x$ 到 $1/x$。
积分 $I = I_1 + I_2$。
$I_1 = \int_0^{1/2} dx \int_x^{4x} (x+2y) dy = \int_0^{1/2} [xy + y^2]_x^{4x} dx = \int_0^{1/2} (4x^2+16x^2 - (x^2+x^2)) dx = \int_0^{1/2} 18x^2 dx = [6x^3]_0^{1/2} = 6/8 = 3/4$。
$I_2 = \int_{1/2}^1 dx \int_x^{1/x} (x+2y) dy = \int_{1/2}^1 [xy+y^2]_x^{1/x} dx = \int_{1/2}^1 (1 + 1/x^2 - (x^2+x^2)) dx = \int_{1/2}^1 (1 + x^{-2} - 2x^2) dx$。
$= [x - 1/x - \frac{2}{3}x^3]_{1/2}^1 = (1 - 1 - 2/3) - (1/2 - 2 - \frac{2}{3}(1/8)) = -2/3 - (-1.5 - 1/12) = -2/3 + 1.5 + 1/12$。
$= -8/12 + 18/12 + 1/12 = 11/12$。
$I = 3/4 + 11/12 = 9/12 + 11/12 = 20/12 = 5/3$。
\end{enumerate}

\item \textbf{最终答案}
\textbf{答案}：$ 5/3 $
\end{detailedanalysis}
\end{redenv}

% ========== 第11题 ==========
\problemrule
\section*{【题目11】X-型区域积分}
求 $ \iint_{D} ye^{xy} dx dy $ ，其中 $D$ 是由 $ xy = 1, x = 2, y = 1 $ 所围成的区域。

\begin{redenv}
\analysisrule
\subsection*{逐步解析}

\begin{detailedanalysis}
\item \textbf{详细求解过程}
\begin{enumerate}[label=(\arabic*)]
\item \textbf{确定区域}：
交点：$xy=1, y=1 \implies x=1$。
$xy=1, x=2 \implies y=1/2$。
$x=2, y=1 \implies (2,1)$。
区域 $x \in [1, 2]$。
下界 $y=1/x$，上界 $y=1$。
\item \textbf{积分计算}：
先积 $y$ 还是 $x$？
被积函数 $ye^{xy}$，对 $x$ 积分容易：$\int ye^{xy} dx = e^{xy}$。
所以先 $x$ 后 $y$？
看区域：$y \in [1/2, 1]$。
左界 $x=1/y$，右界 $x=2$。
$$ I = \int_{1/2}^1 dy \int_{1/y}^2 y e^{xy} dx = \int_{1/2}^1 dy [e^{xy}]_{1/y}^2 = \int_{1/2}^1 (e^{2y} - e^1) dy $$
$$ = [\frac{1}{2}e^{2y} - ey]_{1/2}^1 = (\frac{1}{2}e^2 - e) - (\frac{1}{2}e - \frac{1}{2}e) = \frac{1}{2}e^2 - e $$
\end{enumerate}

\item \textbf{最终答案}
\textbf{答案}：$ \frac{1}{2}e^2 - e $
\end{detailedanalysis}
\end{redenv}

% ========== 第12题 ==========
\problemrule
\section*{【题目12】直角坐标积分}
求二重积分 $ \iint_{D}\frac{y}{x}dxdy $ ，其中 $D$ 是由 $ x=e,y=\ln x,y=0 $ 所围成的区域。

\begin{redenv}
\analysisrule
\subsection*{逐步解析}

\begin{detailedanalysis}
\item \textbf{详细求解过程}
\begin{enumerate}[label=(\arabic*)]
\item \textbf{确定区域}：
$y=0, x=e$。
曲线 $y=\ln x$ 过 $(1,0)$ 和 $(e, 1)$。
区域 $x \in [1, e], 0 \le y \le \ln x$。
\item \textbf{积分计算}：
$$ I = \int_1^e \frac{dx}{x} \int_0^{\ln x} y dy = \int_1^e \frac{1}{x} [\frac{1}{2}y^2]_0^{\ln x} dx $$
$$ = \frac{1}{2} \int_1^e \frac{(\ln x)^2}{x} dx $$
令 $u=\ln x, du=dx/x$。
$$ = \frac{1}{2} \int_0^1 u^2 du = \frac{1}{2} [\frac{1}{3}u^3]_0^1 = \frac{1}{6} $$
\end{enumerate}

\item \textbf{最终答案}
\textbf{答案}：$ \frac{1}{6} $
\end{detailedanalysis}
\end{redenv}

% ========== 第13题 ==========
\problemrule
\section*{【题目13】直角坐标积分}
计算 $ \iint_{D} \frac{x}{y} d\sigma $ ，其中 $D$ 是由 $ y = 1, y = x^{2}, x = 2 $ 所围成的闭区域。

\begin{redenv}
\analysisrule
\subsection*{逐步解析}

\begin{detailedanalysis}
\item \textbf{详细求解过程}
\begin{enumerate}[label=(\arabic*)]
\item \textbf{确定区域}：
$y=x^2$ 与 $y=1$ 交于 $(1,1)$。
$x=2$ 时 $y=4$。
区域 $x \in [1, 2]$，$1 \le y \le x^2$。
\item \textbf{积分计算}：
$$ I = \int_1^2 x dx \int_1^{x^2} \frac{dy}{y} = \int_1^2 x [\ln y]_1^{x^2} dx $$
$$ = \int_1^2 x \ln(x^2) dx = 2 \int_1^2 x \ln x dx $$
分部积分：
$$ 2 [\frac{x^2}{2}\ln x - \int \frac{x^2}{2}\frac{1}{x} dx]_1^2 = 2 [ \frac{x^2}{2}\ln x - \frac{x^2}{4} ]_1^2 $$
$$ = [x^2\ln x - \frac{x^2}{2}]_1^2 = (4\ln 2 - 2) - (0 - 1/2) = 4\ln 2 - 1.5 $$
\end{enumerate}

\item \textbf{最终答案}
\textbf{答案}：$ 4\ln 2 - \frac{3}{2} $
\end{detailedanalysis}
\end{redenv}

% ========== 第14题 ==========
\problemrule
\section*{【题目14】直角坐标积分}
计算二重积分 $ \iint_{D} x\sqrt{y} $ 其中 $D$ 是由 $ y = \sqrt{x}, y = x^{2} $ 围成的区域。

\begin{redenv}
\analysisrule
\subsection*{逐步解析}

\begin{detailedanalysis}
\item \textbf{详细求解过程}
\begin{enumerate}[label=(\arabic*)]
\item \textbf{确定区域}：
$x^2 = \sqrt{x} \implies x^4 = x \implies x=0, 1$。
区域 $x \in [0, 1]$，$x^2 \le y \le \sqrt{x}$。
\item \textbf{积分计算}：
$$ I = \int_0^1 x dx \int_{x^2}^{\sqrt{x}} y^{1/2} dy = \int_0^1 x [\frac{2}{3}y^{3/2}]_{x^2}^{\sqrt{x}} dx $$
$$ = \frac{2}{3} \int_0^1 x ( (\sqrt{x})^{3/2} - (x^2)^{3/2} ) dx = \frac{2}{3} \int_0^1 x (x^{3/4} - x^3) dx $$
$$ = \frac{2}{3} \int_0^1 (x^{7/4} - x^4) dx = \frac{2}{3} [\frac{4}{11}x^{11/4} - \frac{1}{5}x^5]_0^1 $$
$$ = \frac{2}{3} (\frac{4}{11} - \frac{1}{5}) = \frac{2}{3} \frac{20-11}{55} = \frac{2}{3} \frac{9}{55} = \frac{6}{55} $$
\end{enumerate}

\item \textbf{最终答案}
\textbf{答案}：$ \frac{6}{55} $
\end{detailedanalysis}
\end{redenv}

\end{document}
